\subsection{Positioning Summary}

The preceding literature addresses complementary parts of the diagnostic problem, but usually under different task definitions and evidence assumptions. UAV fault-detection studies and multivariate time-series detectors primarily determine whether telemetry is abnormal, while AeroTSBoost provides the closest real-flight, temporal-statistical detector reference for the present setting \cite{keipour2019automatic,wei2026aerotsboost}. Explainability research supplies attribution and faithfulness tests, but its outputs remain tied to model inputs unless an explicit aggregation and architecture mapping is defined \cite{lundberg2020trees,hooker2019benchmark}. Conversely, spectrum-based software fault localization ranks program entities from test executions and known pass/fail outcomes, evidence that routine real-flight logs do not provide \cite{wong2016survey}. These research lines are therefore individually relevant, yet none of their typical formulations spans the complete path from a detected telemetry window to validated diagnostic evidence and version-aware PX4 software-module suspects.

HS-AeroTS is positioned as an auditable integration of these evidence levels rather than as a replacement for their underlying methods. It couples real-flight runtime anomaly detection with four-domain diagnosis, evaluates the resulting hierarchical SHAP evidence through conservation, masking, and label-consistency tests, and propagates topic evidence over a uORB graph reconstructed at the selected majority firmware commit. Its evaluation also reports fixed chronological, purged, and leave-log-out protocols, together with external transfer and controlled replay. Crucially, conditional module-ranking results are reported alongside the detector-gated replay failure, preventing onset-conditioned success from being interpreted as reliable end-to-end localization. Table~\ref{tab:related_work_positioning} summarizes this distinction; the entries describe the typical scope of each research line, not comparative predictive performance.

\begin{table}[H]
\caption{Positioning of representative research lines and HS-AeroTS across the four evidence dimensions emphasized in this study. ``Task-dependent'' indicates that the method can be attached to different predictive tasks but does not itself define the listed diagnosis level.}
\label{tab:related_work_positioning}
\centering
\footnotesize
\renewcommand{\arraystretch}{1.18}
\begin{tabularx}{\textwidth}{>{\RaggedRight\arraybackslash}p{3.0cm} >{\RaggedRight\arraybackslash}p{2.35cm} >{\RaggedRight\arraybackslash}p{2.25cm} >{\RaggedRight\arraybackslash}p{2.55cm} >{\RaggedRight\arraybackslash}X}
\toprule
\textbf{Research line} & \textbf{Runtime anomaly detection} & \textbf{Fault-domain diagnosis} & \textbf{Explanation validation} & \textbf{Software-module evidence} \\
\midrule
UAV fault detection and diagnosis \cite{fourlas2021survey,puchalski2022uav,keipour2021alfa} & Flight- or signal-level detection & Fault or failure categories & Not a primary evaluation target & No source-architecture mapping \\
General multivariate time-series anomaly detection \cite{su2019omnianomaly,deng2021gdn,tuli2022tranad} & Anomaly scores or labels & Not intrinsic & Usually outside detector evaluation & None \\
AeroTSBoost \cite{wei2026aerotsboost} & Real-flight binary detection & Not addressed & Not addressed & None \\
Attribution and explanation evaluation \cite{bento2021timeshap,adebayo2018sanity,hooker2019benchmark} & Task-dependent & Task-dependent & Attribution sanity or perturbation tests & No PX4 architecture mapping \\
Spectrum-based software fault localization \cite{jones2002visualization,wong2016survey,pearson2017evaluating} & No telemetry detector & No physical fault domain & Ranking against known faulty entities & Execution-spectrum entity ranking \\
\textbf{HS-AeroTS} & Real-flight binary detection & Four UAV-SEAD domains & Conservation, masking, and semantic-map agreement & Matching-commit uORB module suspects \\
\bottomrule
\end{tabularx}
\end{table}
